\chapter{The String Interner}
\label{ch:interner}

Entity identifiers in ingested logs are long and massively
repetitive --- a single \code{SecurityEvent} table in Sentinel will
mention the same host name hundreds of thousands of times. Storing
those strings as owned \code{String} inside \code{CompactRelation}
would balloon the edge list. The engine side-steps this with a
bidirectional interner that maps each unique string to a 32-bit
handle \StrId{}.

\section{Structure}

\begin{figure}[h]
\centering
\begin{lstlisting}[style=rust]
pub struct StrId(u32);

pub struct StringInterner {
    map:     HashMap<Arc<str>, StrId>,  // reverse lookup
    strings: Vec<Arc<str>>,             // forward lookup by index
}
\end{lstlisting}
\end{figure}

Both containers hold \code{Arc<str>} pointing at the same
heap-allocated slice, so each intern costs exactly one allocation
(the \code{Arc::from(\&str)} inside the \code{intern} path) plus two
refcount bumps. The alternative --- \code{String} in the map, plain
\code{String} or a second copy in the vector --- would double the
allocator pressure.

\begin{invariant}[Intern bijection]
For every index $i \in [0, \code{strings.len()})$ we have
\code{map[strings[i]] == StrId(i as u32)}. The two containers are
populated atomically inside \code{intern}; a panic during insertion
cannot leave them out of sync because a half-inserted entry is
not visible to subsequent lookups (the \code{map.insert} is the
last action on the happy path).
\end{invariant}

\section{Algorithms}

\begin{algorithm}
\caption{\textsc{Intern}}\label{alg:intern}
\KwIn{interner $I$, string $s$}
\KwOut{\StrId{} assigned to $s$}
\If{$s \in I.\text{map}$}{\Return $I.\text{map}[s]$\;}
$\text{id} \gets I.\text{strings.len()}$ as \code{u32}\;
$\text{arc} \gets \code{Arc::from}(s)$\Comment*[r]{1 heap alloc}
\code{I.strings.push(arc.clone())}\Comment*[r]{refcount++ }
\code{I.map.insert(arc, StrId(id))}\Comment*[r]{refcount++ }
\Return \code{StrId(id)}\;
\end{algorithm}

\begin{algorithm}
\caption{\textsc{Resolve}}\label{alg:resolve}
\KwIn{interner $I$, \StrId{} $x$}
\KwOut{$\&$\code{str} for $x$, or $\bot$ if out-of-range}
\If{$x.0 \geq I.\text{strings.len()}$}{\Return $\bot$\;}
\Return $\&* I.\text{strings}[x.0]$\;
\end{algorithm}

\begin{complexity}
\textsc{Intern} is $\Ord{1}$ amortized; the \code{HashMap} uses
ahash which hits very few collisions on real entity IDs, and the
\code{Vec::push} amortises. The heap allocation occurs on the miss
path only. \textsc{Resolve} is $\Thet{1}$: a bounds check and a
pointer deref. Space consumption is
\[
    \Thet{\sum_{s \in \mathcal{U}} |s|} + \Thet{|\mathcal{U}|}
\]
where $\mathcal{U}$ is the set of unique strings --- the first
term is the heap content, the second is the $\Ord{1}$-bytes-per-entry
overhead of the two containers.
\end{complexity}

\section{Metrics}

The engine exposes interner memory usage via
\code{StringInterner::metrics()}:

\begin{lstlisting}[style=rust]
pub struct InternerMetrics {
    pub unique_strings: usize,
    pub approx_bytes:   usize,   // sum of Arc<str>.len()
}
\end{lstlisting}

These are the trigger metrics for the deferred \emph{generational}
or \emph{partitioned} back-ends (Chapter~\ref{ch:interner-trait}):
when \code{unique\_strings} crosses $\sim$$10^7$, or when a
long-running streaming session's \code{approx\_bytes} grows
monotonically past a few hundred MB, the operator knows to switch.

\section{The trait seam}

The P2-F refactor introduced \code{StringInternerBackend} as the
swap-in point:

\begin{lstlisting}[style=rust]
pub trait StringInternerBackend {
    fn intern(&mut self, s: &str) -> StrId;
    fn resolve(&self, id: StrId) -> &str;
    fn try_resolve(&self, id: StrId) -> Option<&str>;
    fn get(&self, s: &str) -> Option<StrId>;
    fn len(&self) -> usize;
    fn is_empty(&self) -> bool;
    fn reserve(&mut self, additional: usize);
    fn metrics(&self) -> InternerMetrics;
}
\end{lstlisting}

\code{StringInterner} is the only impl today. When the engine
acquires evidence that the global-interner assumption breaks
(multi-tenant process, 24$\times$7 streaming past $10^7$ unique
strings), a \code{GenerationalInterner} or
\code{PartitionedInterner} drops behind the same trait; the 121+
existing call sites keep compiling because they hold a concrete
\code{StringInterner} and the inherent \code{metrics()} method
delegates to the trait body.

\begin{inthecode}
\filepath{graph\_hunter\_core/src/interner.rs:1} --- struct + trait
+ metrics.\\
\filepath{graph\_hunter\_core/src/interner.rs:115} --- unit tests
pinning the bijection.
\end{inthecode}
